% META: {"id": "1NSI-RES-AM-EXTRAIT", "chapitre": "1NSI-RESEAUX", "type_objet": "amenagee", "capacites": ["P-ARCH-02A", "P-ARCH-02B", "P-ARCH-04A"], "capacites_codes": ["C1", "C2", "C4"], "derive_de": ["1NSI-RES-EX-001", "1NSI-RES-EX-002", "1NSI-RES-EX-004"], "profil_amenagement": ["SEQUENCAGE", "ALLEGEMENT_ENONCE", "REPONSE_FERMEE", "AERATION", "AMORCE_FOURNIE", "RETRAIT_DOUBLE_TACHE"], "mode_creation": "adaptation", "sources_inspiration": [], "status": "generated"}
\section*{Version aménagée --- Extrait Réseaux}
\textit{Consignes séquencées, mise en page aérée, énoncés allégés.}

\bigskip

\subsection*{Exercice 1 --- Découper un message en paquets (C1)}

Le message \texttt{PROTOCOLERESEAU} est découpé en paquets de $4$ caractères.

\bigskip

\textbf{Étape 1.} Compte les caractères du message.

\medskip
Nombre de caractères : \dotfill

\bigskip

\textbf{Étape 2.} Découpe. Le premier paquet est donné.

\bigskip

\begin{center}
\begin{tabular}{|c|c|}
\hline
\textbf{Numéro} & \textbf{Contenu} \\
\hline
$1$ & \texttt{PROT} \\
\hline
$2$ & \dotfill \\
\hline
$3$ & \dotfill \\
\hline
$4$ & \dotfill \\
\hline
\end{tabular}
\end{center}

\bigskip

\textbf{Étape 3.} Le dernier paquet est-il plein ?

\medskip
\begin{center}
\fbox{Oui} \qquad \fbox{Non}
\end{center}

\bigskip

\textbf{Étape 4.} Pourquoi chaque paquet porte-t-il un numéro ? Entoure.

\medskip
\begin{center}
\fbox{pour les remettre dans l'ordre à l'arrivée} \qquad \fbox{pour aller plus vite}
\end{center}

\bigskip
\bigskip

\subsection*{Exercice 2 --- Le protocole du bit alterné (C2)}

On envoie \texttt{X}, puis \texttt{Y}, puis \texttt{Z}. L'accusé du \textbf{premier}
envoi se perd.

\bigskip

\textbf{Étape 1.} Complète le tableau. Une ligne à la fois.

\bigskip

\begin{center}
\begin{tabular}{|c|l|c|c|}
\hline
\textbf{Envoi} & \textbf{Ce qui se passe} & \textbf{Bit} & \textbf{Reçu ?} \\
\hline
$1$ & envoi de \texttt{X} & $0$ & oui \\
\hline
$2$ & accusé perdu, on renvoie \texttt{X} & $0$ & \dotfill \\
\hline
$3$ & envoi de \texttt{Y} & \dotfill & \dotfill \\
\hline
$4$ & envoi de \texttt{Z} & \dotfill & \dotfill \\
\hline
\end{tabular}
\end{center}

\bigskip

\textbf{Étape 2.} À l'envoi $2$, le récepteur reçoit \texttt{X} une seconde fois.
Que fait-il ?

\medskip
\begin{center}
\fbox{il le garde deux fois} \qquad \fbox{il le jette, mais renvoie un accusé}
\end{center}

\bigskip

\textbf{Étape 3.} Combien de messages différents le récepteur a-t-il gardés ?

\medskip
\dotfill

\bigskip
\bigskip

\subsection*{Exercice 3 --- Capteur ou actionneur ? (C4)}

Un lampadaire s'allume quand la luminosité descend sous $300$ lux.

\bigskip

\textbf{Étape 1.} Classe les deux éléments.

\bigskip

\begin{center}
\begin{tabular}{|l|c|}
\hline
\textbf{Élément} & \textbf{Capteur ou actionneur ?} \\
\hline
La cellule qui mesure la luminosité & \dotfill \\
\hline
La lampe qui s'allume & \dotfill \\
\hline
\end{tabular}
\end{center}

\medskip
\textit{Aide :} un capteur \textbf{mesure}, un actionneur \textbf{agit}.

\bigskip

\textbf{Étape 2.} Complète la condition. Le début est donné.

\begin{python}
if luminosite ______ 300:
    allumer_lampe()
\end{python}

\medskip
Le signe manquant : \fbox{\texttt{<}} \qquad ou \qquad \fbox{\texttt{>}}
